% ======================================================================
% preprints/jphq_06_structural_tension_nonoptimizable/sections/07_falsifiability.tex
% Journal-grade content (100/100). ASCII-only.
% ======================================================================

\section{Falsifiability and Conditions of Refutation}

This section specifies how the central claim of this preprint may be
falsified. Falsifiability is defined strictly in terms of formal
consistency with the Ontology of Continua (OC) and
\texttt{ETS.K\_EA.v1.1}. No empirical validation claims are made.

\subsection{Target Claim}

The sole claim under examination is the following:

\begin{quote}
Structural tension $\TEA$ in the enterprise continuum $\KEA$ is not an
optimizable quantity. It cannot be minimized, maximized, or treated as a
scalar objective without violating admissibility, phase logic, or
continuity semantics as defined in OC/ETS.
\end{quote}

\subsection{Admissible Modes of Refutation}

The claim is refuted if \emph{any one} of the following can be
demonstrated \emph{without modifying} OC Core or
\texttt{ETS.K\_EA.v1.1}:

\begin{enumerate}
  \item A scalar objective function
        \[
        J : \OmegaEA \rightarrow \mathbb{R}
        \]
        is constructed that represents $\TEA$ and whose minimization or
        maximization preserves admissibility of $\OmegaEA$ across all
        threshold components.
  \item Structural tension $\TEA$ is shown to admit a total or partial
        order on $\OmegaEA$ that is fully compatible with phase
        precedence (\PHASESTOP\ $\succ$ \PHASECOLLAPSE\ $\succ$
        \PHASEINERTIA\ $\succ$ \PHASESUCCESS) and does not refine phases
        internally.
  \item A continuous deformation or parametrization of $\TEA$ exists
        that preserves logical phase boundaries while enabling
        gradient-based or extremal optimization within $\OmegaEA$.
  \item Optimization of $\TEA$ can be shown to preserve non-zero
        continuity $\kEA$ under all admissible transitions, without
        invoking metric substitution or additional primitives.
\end{enumerate}

Satisfaction of any single condition suffices to falsify the present
result.

\subsection{Explicit Non-Falsifiers}

The following do \emph{not} constitute refutation of the claim:

\begin{itemize}
  \item Empirical case studies reporting improved outcomes correlated
        with reduced stress, pressure, or instability.
  \item Heuristic or ordinal scoring schemes that label threshold
        violations as ``high'' or ``low'' tension.
  \item Simulations that reinterpret thresholds as soft constraints,
        penalties, or barrier terms.
  \item Reinterpretations of $\TEA$ as a proxy for risk, cost, loss, or
        utility.
\end{itemize}

All such approaches modify the semantics of structural tension and
therefore fall outside the formal object addressed by this paper.

\subsection{Internal Consistency Refutation}

The claim is also refuted if an internal inconsistency is identified
within OC/ETS itself, including but not limited to:
(i) contradiction between threshold definitions and phase logic;
(ii) existence of admissible cycles with $f_k > 0$ for any required
threshold component; or
(iii) coexistence of optimization semantics with admissible STOP states.

Such inconsistencies would invalidate the impossibility results by
undermining their formal premises.

\subsection{Scope of Falsification}

Refutation applies only to the specific formal statement addressed here.
It does not invalidate:
(i) alternative enterprise models,
(ii) optimization-based frameworks outside OC/ETS, or
(iii) empirical or pragmatic methods unconstrained by admissibility
logic.

The falsifiability conditions are therefore narrow by construction and
coextensive with the scope of the claim.

%%% End of section %%%
